\section{Contrôle du repo à partir des contrôles forward et dividendes}\label{sec-control}

\subsection{Lignes de contrôle}
On regroupe les taux de repo zéro-coupon aux piliers et les niveaux de dividendes par période dans
$\theta=(b,L)\in\R^K\times\R^Y$. Le spot, la courbe des taux et le profil d'allocation sont fixés, de sorte que les
forwards aux piliers forment une application régulière $F(\theta)\in\R^K$. D'après la \cref{prop-bootstrap},
l'\emph{application repo implicite} $\beta(F,L)$, l'unique $b$ tel que $F(b,L)=F$, est bien définie.

Les marquages FO sont $\theta^\FO=(b^\FO,L^\FO)$. Totem fournit des forwards $F^\T$ et des niveaux de dividendes $L^\T$,
d'où $b^\T=\beta(F^\T,L^\T)$. Une position a une valeur $V(\theta)$, avec des sensibilités FO
$s_b=\partial_bV(\theta^\FO)\in\R^K$ et $s_D=\partial_LV(\theta^\FO)\in\R^Y$. On note $\Delta L=L^\T-L^\FO$,
$\Delta b=b^\T-b^\FO$ et $\Delta b^F=\beta(F^\T,L^\FO)-b^\FO$ : l'écart de repo implicité par les forwards Totem
\emph{avec les dividendes FO}.

\begin{definition}[Lignes de contrôle]\label{def-lines}
\[
\text{ligne dividende}=s_D\cdot\Delta L,\qquad\text{ligne repo}=s_b\cdot\Delta b,\qquad\text{ligne forward}=s_b\cdot\Delta b^F,
\]
et le \emph{total} est la somme des lignes dividende et repo.
\end{definition}

\begin{definition}[Exposition linéaire en forward]\label{def-fwdlinear}
Une position est linéaire en forward si $V(\theta)=\Phi\bigl(F(\theta)\bigr)$ pour une fonction $\Phi$ de classe $C^2$ :
sa valeur ne dépend du repo et des dividendes qu'à travers les forwards (options européennes avec une nappe de
volatilité qui n'en dépend pas, forwards, futures).
\end{definition}

\subsection{La jacobienne du repo implicite}
\begin{lemma}\label{lem-jacobian}
$J:=\partial_L\beta(F,L)=-(\partial_bF)^{-1}\partial_LF$. La matrice $\partial_bF$ est triangulaire inférieure avec des
termes diagonaux $\partial F_k/\partial b_k\le-T_kF_k<0$, donc $J$ existe. Au premier ordre,
\begin{equation}\label{eq-jacobian}
J_{ky}\approx-\frac{A_{ky}}{T_k\,(S-D_k)}.
\end{equation}
\end{lemma}
\begin{proof}
On différentie $F(\beta(F,L),L)=F$ par rapport à $L$ : $\partial_bF\,J+\partial_LF=0$. Avec un repo forward constant
entre piliers, $F_k$ ne dépend que de $b_1,\dots,b_k$, donc $\partial_bF$ est triangulaire inférieure ; sa diagonale vaut
$T_k\,\Phi_k'(B_k)\le-T_kF_k$ d'après le \cref{lem-monotone}. Pour le premier ordre, on fixe les forwards dans
\cref{eq-additivity} : $\delta B_k=-\delta D_k/(S-D_k)=-A_{ky}\delta L_y/(S-D_k)$, et on divise par $T_k$.
\end{proof}

\subsection{L'identité de contrôle}
\begin{theorem}[Total indépendant des dividendes]\label{thm-control}
Pour toute position,
\begin{equation}\label{eq-control}
\text{ligne dividende}+\text{ligne repo}=\text{ligne forward}+g\cdot\Delta L+O(\|\Delta\|^2),\qquad g=s_D+s_b\,J,
\end{equation}
avec $J$ évaluée en $(F^\T,L^\FO)$ et $\|\Delta\|=\|\Delta L\|+\|\Delta b^F\|$. Si la position est linéaire en forward,
alors $g=O(\|\Delta\|)$ : le total égale la ligne forward au second ordre près et ne dépend donc pas du panel
dividendes au premier ordre.
\end{theorem}
\begin{proof}
Un développement au premier ordre de $\beta$ en $L$ donne $\beta(F^\T,L^\T)=\beta(F^\T,L^\FO)+J\Delta L+O(\|\Delta L\|^2)$,
soit $\Delta b=\Delta b^F+J\Delta L+O(\|\Delta L\|^2)$. D'où
$s_D\cdot\Delta L+s_b\cdot\Delta b=s_b\cdot\Delta b^F+(s_D+s_bJ)\cdot\Delta L+O(\|\Delta\|^2)$.
Pour $V=\Phi\circ F$, $s_b=\nabla\Phi\,\partial_bF(\theta^\FO)$ et $s_D=\nabla\Phi\,\partial_LF(\theta^\FO)$, donc
$g=\nabla\Phi\,\bigl[\partial_LF+\partial_bF\,J\bigr](\theta^\FO)$. En $\theta'=(\beta(F^\T,L^\FO),L^\FO)$, le crochet est nul
d'après le \cref{lem-jacobian} ; comme $\theta'-\theta^\FO=(\Delta b^F,0)$ et que les dérivées sont lipschitziennes, le
crochet en $\theta^\FO$ est un $O(\|\Delta b^F\|)$, donc $g\cdot\Delta L=O(\|\Delta\|^2)$.
\end{proof}

\begin{theorem}[Exactitude de la réévaluation complète]\label{thm-exact}
Si la position est linéaire en forward et que toutes ses échéances sont des piliers Totem, le contrôle en
réévaluation complète $V(\theta^\T)-V(\theta^\FO)=\Phi(F^\T)-\Phi\bigl(F(\theta^\FO)\bigr)$ ne dépend pas du tout des
dividendes Totem $L^\T$.
\end{theorem}
\begin{proof}
Par construction du repo implicite, $F(\theta^\T)=F\bigl(\beta(F^\T,L^\T),L^\T\bigr)=F^\T$ quel que soit $L^\T$.
\end{proof}

\begin{remark}[Échéances entre piliers]
Pour une échéance située entre deux piliers, le forward dépend du portage interpolé et des dividendes alloués entre
les deux : une dépendance résiduelle à $L^\T$ subsiste. Elle est de l'ordre des dividendes détachés entre l'échéance et
les piliers voisins, c'est-à-dire exactement les effets de date et de profil des \cref{prop-timing,prop-profile}, que
le détecteur surveille.
\end{remark}

\begin{proposition}[Terme du second ordre]\label{prop-second}
Pour un pilier, avec $F(b,D)=e^{R-bT}(S-D)$ et $V=\Phi\circ F$, le total au premier ordre vaut
\begin{equation}\label{eq-second}
\Phi'(F)\,\Delta F-\Phi'(F)\,F\Bigl[\tfrac12T^2\Delta b^2+\frac{T\,\Delta b\,\Delta D}{S-D}\Bigr]+O(\|\Delta\|^3),
\end{equation}
alors que le contrôle exact vaut $\Phi'(F)\Delta F+\tfrac12\Phi''(F)\Delta F^2+O(\|\Delta\|^3)$. La dépendance du total au
premier ordre vis-à-vis du panel dividendes est donc entièrement contenue dans le crochet.
\end{proposition}
\begin{proof}
On développe $F$ au second ordre : $F_b=-TF$, $F_D=-F/(S-D)$, $F_{bb}=T^2F$, $F_{bD}=TF/(S-D)$, $F_{DD}=0$, d'où
$\Delta F=F_b\Delta b+F_D\Delta D+\tfrac12T^2F\Delta b^2+TF\Delta b\,\Delta D/(S-D)+O(3)$. Le total au premier ordre vaut
$\Phi'(F_b\Delta b+F_D\Delta D)$ ; on substitue.
\end{proof}

\subsection{Ce que la répartition coûte à un portefeuille}
\begin{proposition}[Matérialité de la répartition]\label{prop-materiality}
Supposons que l'erreur de consensus dividendes $\Delta L$ ait une covariance $\Sigma_L$ (bruit de consensus), les forwards
Totem étant fixés. Au premier ordre, la variance du total vaut $g^\top\Sigma_Lg$, contre $s_D^\top\Sigma_Ls_D$ pour la
seule ligne dividende. La part de l'incertitude de la ligne dividende compensée par la ligne repo vaut
$1-\sqrt{g^\top\Sigma_Lg\,/\,s_D^\top\Sigma_Ls_D}$.
\end{proposition}
\begin{proof}
D'après le \cref{thm-control}, le total est la ligne forward, qui ne fait pas intervenir $L^\T$, plus $g\cdot\Delta L$.
\end{proof}

\begin{corollary}[Ratio de compensation]
Pour un portefeuille linéaire en forward, un mouvement du panel dividendes produit des lignes dividende et repo
importantes et opposées, avec un total faible. Un ratio $1-|\text{total}|/(|\text{ligne div.}|+|\text{ligne repo}|)$
proche de un signale donc un artefact de répartition, pas une erreur de marquage.
\end{corollary}

$g$ se calcule à partir des seules sensibilités FO et de la jacobienne du bootstrap ; aucune revalorisation n'est
nécessaire.

\subsection{Exemples chiffrés}
\begin{example}[Un call]\label{ex-call}
Un call à la monnaie à un an (spot $100$, taux $3\,\%$, volatilité $25\,\%$, dividende FO $3{,}5$, repo FO $60$\,pb),
avec un forward Totem inférieur de $20$\,pb à celui du FO et trois panels dividendes différents :
\begin{center}\input{generated/\lang/tab_cancel}\end{center}
Le repo implicite varie de $125$\,pb d'un panel à l'autre et chaque ligne d'environ $0{,}3$, mais le total ne bouge pas.
\end{example}

\begin{example}[Trois portefeuilles]\label{ex-books}
Le pricer du prototype valorise trois portefeuilles sur onze piliers de repo et six périodes de dividendes : A (calls,
puts et un forward), B (A plus un swap de dividendes 2028 vendu) et C (A plus une commission de prêt de titres à un
an). Les forwards Totem sont inférieurs de $20$\,pb à ceux du FO et le panel dividendes Totem déplace les périodes
2027 et 2028 de $-0{,}4$, $0$ ou $+0{,}5$ :
\begin{center}\input{generated/\lang/tab_books}\end{center}
Pour le portefeuille A, le total reste à \vBookATotal{} alors que chaque ligne varie de plus de six unités. Pour B et C,
le total bouge, exactement de la ligne non-forward $g\cdot\Delta L$ : le swap de dividendes est une véritable
exposition aux dividendes et la commission de prêt une véritable exposition au repo.
\end{example}

\begin{figure}[ht]
\centering\includegraphics[width=\linewidth]{figures/\lang/cancellation.pdf}
\caption{Ligne dividende, ligne repo et leur total quand le panel dividendes Totem se déplace (portefeuille A à
gauche, B à droite). Points : réévaluation complète. Pour une exposition linéaire en forward, le total est plat
(\cref{thm-control,thm-exact}) ; le swap de dividendes ajoute une véritable exposition aux dividendes dont la ligne est
exactement $g\cdot\Delta L$.}
\label{fig-cancellation}
\end{figure}

\begin{example}[Matérialité]\label{ex-materiality}
Avec une dispersion de consensus dividendes de $5\,\%$ de chaque période, dix contributeurs et une corrélation
$e^{-|\Delta y|/2}$ entre périodes, la \cref{prop-materiality} donne
\begin{center}\input{generated/\lang/tab_materiality}\end{center}
\vMatA{} de l'incertitude de la ligne dividende se compense pour le portefeuille A ; \vMatB{} et \vMatC{} pour B et C.
\end{example}

\subsection{Quand la répartition compte}
L'exposition non-forward $g$ est non nulle pour les swaps et futures de dividendes (dividende pur), pour l'emprunt de
titres et les jambes de financement des swaps de rendement total (repo pur), pour l'exercice américain (la frontière
d'exercice anticipé dépend séparément des dividendes et du repo) et pour les exotiques valorisés avec une dynamique de
dividendes discrets (le modèle de dividendes modifie la volatilité du spot). Pour ces portefeuilles, la répartition
doit être contrôlée directement, et la qualité de la courbe de repo implicite (\cref{sec-detection}) compte.

\begin{insight}
\begin{itemize}[leftmargin=12pt]
\item Contrôler la \emph{ligne forward} et la \emph{ligne non-forward} : la première n'a besoin que des forwards Totem
et des dividendes FO ; la seconde n'a besoin des dividendes Totem que là où $g\neq0$.
\item Des lignes dividende et repo importantes qui se compensent, avec un total faible, sont un artefact de répartition.
\item En valorisation prudente~\cite{eu2016}, si les incertitudes sont agrégées par paramètre de valorisation sans
compensation, un artefact de répartition est compté deux fois ; définir le paramètre comme le forward là où $g\approx0$
supprime ce double comptage.
\end{itemize}
\end{insight}
